\documentclass[11pt]{article}
\usepackage[margin=1in]{geometry}
\usepackage{amsmath,amssymb}
\usepackage{enumitem}
\usepackage{xcolor}
\usepackage{parskip}
\usepackage{booktabs}
\usepackage{array}
\usepackage{hyperref}
\hypersetup{colorlinks=true, urlcolor=blue!45!black, linkcolor=blue!45!black}
\title{\vspace{-2em}\textbf{Spatial Genomics / Spatial Transcriptomics}\\[0.3em]\large Keeping the map: measuring expression with position retained\\[0.6em]\normalsize\textit{Genomics Hackathon --- methods primer}}
\author{}\date{}

\begin{document}
\maketitle
\tableofcontents

\bigskip
\fbox{\parbox{\linewidth}{\small
\textbf{Sourcing note.} This is a synthesis primer for the hackathon. Single-cell and sequencing foundations are
grounded in \emph{Genomics Lecture~1}. The spatial-transcriptomics content is compiled from standard methods
knowledge plus the project's own framing, because it is \emph{not} present in the typed English lecture notes ---
those notes contain no dedicated spatial-transcriptomics lecture (the lecture called ``spatial,'' Lecture~12, is
about Hi-C / 3D-genome folding, which is a different meaning of ``spatial''). The single-cell recap below follows
Lecture~1.
}}
\bigskip

\noindent\textit{From dissociated single cells to expression with a street address --- and using a spatial map as a
ruler to decode disease.} Builds on Genomics Lecture~1 (sequencing foundations); audience: students who know
scRNA-seq basics.

\section{Where this material comes from}

\subsection{Recap: how we turn a tissue into numbers (from Lecture 1)}
Sequencing reads molecules in parallel. Sanger uses chain-terminating ddNTPs that stop synthesis randomly across
many template copies, then sorts fragments by length. Illumina spreads fragments onto a flow cell, amplifies each
into a local cluster, and photographs one base per cycle (``sequencing by synthesis''). Crucially, the physical
position of a cluster on the flow cell is meaningless --- it just says ``a read is here''; genomic identity comes
only from alignment.

Hold onto that last point. In ordinary sequencing the $(x,y)$ on the machine is throwaway. The whole idea of
spatial transcriptomics is to make a position $(x,y)$ mean something biological again --- not on the flow cell, but
in the tissue.

\subsection{Single-cell RNA-seq, in one paragraph}
RNA-seq measures gene expression by sequencing the mRNA in a sample and counting reads per gene. Single-cell
RNA-seq (scRNA-seq) does this one cell at a time: the tissue is dissociated into a suspension, each cell is captured
(e.g.\ in a droplet) with a bead carrying a unique \textbf{cell barcode}, every captured mRNA gets a \textbf{UMI}
(unique molecular identifier) so PCR duplicates can be collapsed, and you end up with a cells~$\times$~genes count
matrix. Cluster that matrix and you recover cell types.

\fbox{\parbox{\linewidth}{\small
\textbf{Key vocabulary.} scRNA-seq $=$ single-cell RNA-seq (whole cell). snRNA-seq $=$ single-nucleus RNA-seq ---
only the nucleus is profiled, which is the trick that lets you measure frozen or fibrotic tissue where whole cells
will not survive dissociation. Both are \emph{dissociated} methods.
}}

\section{Why position matters --- what dissociation destroys}
There are two fundamentally different ways to measure a tissue, and the difference is entirely about whether you
keep the cells' addresses.

\subsection{The dissociated view}
To do scRNA-seq / snRNA-seq you must dissociate the tissue --- chop and digest it until cells (or nuclei) float
free in a suspension. The reward is enormous: you read essentially the whole transcriptome of each cell, thousands
of genes, at single-cell resolution. The price is brutal and irreversible:

\fbox{\parbox{\linewidth}{\small
\textbf{What is lost.} The moment cells go into suspension, their position is destroyed. Cell \#5{,}000 and cell
\#5{,}001 in your matrix might have been neighbors, or on opposite sides of the tissue --- you can never tell. You
get a perfect parts list with the assembly diagram thrown away.
}}

\subsection{The spatial view}
Spatial transcriptomics measures expression \emph{without} dissociating --- on an intact tissue section, so each
measurement carries its $(x,y)$ coordinate. You keep the assembly diagram.

\fbox{\parbox{\linewidth}{\small
\textbf{The core distinction.} Dissociated $=$ high transcriptome coverage, position gone. Spatial $=$ expression
tied to a location, usually with some cost in coverage or sensitivity.
}}

\subsection{What only the spatial view can show}
\begin{itemize}[leftmargin=1.4em]
  \item \textbf{Tissue architecture:} map each cell type onto where it actually sits --- layers, structures,
        lesions.
  \item \textbf{Gradients \& axes:} reconstruct how a gene rises or falls along a spatial direction. The textbook
        case is \emph{liver zonation} --- hepatocytes change their metabolic program continuously along the
        porto-central axis (from the portal triad to the central vein). Dissociation flattens this gradient into
        noise; spatial data recovers it as a smooth curve.
  \item \textbf{Niches / neighborhoods:} ask which cell types sit next to which --- e.g.\ immune cells clustering
        around a tumor edge.
  \item \textbf{Spatial ligand--receptor inference:} a signaling cell and a responding cell must be close; only
        with coordinates can you score whether a ligand-expressing cell actually neighbors a receptor-expressing
        one.
\end{itemize}

\section{The spatial methods and their trade-offs}
No method gives you everything. They sit on a triangle of competing goods: \textbf{transcriptome breadth} (how many
genes), \textbf{spatial resolution} (how fine the $(x,y)$), and \textbf{sensitivity} (how many molecules you
actually catch). Two broad families exist.

\paragraph{Sequencing-based (capture \& sequence).}
The tissue is laid on a slide patterned with spatially-barcoded capture spots. mRNA diffuses down, binds a spot, and
inherits that spot's spatial barcode --- the same barcoding trick as scRNA-seq, except the barcode now encodes a
\emph{tissue location}. Whole-transcriptome, but resolution is limited by spot size. (Notice the echo of Lecture~1:
a barcode tells the sequencer where a read belongs --- here ``where'' is a real place.)

\paragraph{Imaging-based (in situ hybridization / sequencing).}
Genes are detected directly inside the intact cell by fluorescent probes imaged under a microscope, often one
molecule at a time. Subcellular resolution, but you can only target a pre-chosen panel of genes, not the whole
transcriptome.

\begin{table}[h]
\centering
\small
\begin{tabular}{>{\raggedright\arraybackslash}p{2.3cm}
                >{\raggedright\arraybackslash}p{2.1cm}
                >{\raggedright\arraybackslash}p{2.5cm}
                >{\raggedright\arraybackslash}p{2.6cm}
                >{\raggedright\arraybackslash}p{4.0cm}}
\toprule
\textbf{Method} & \textbf{Family} & \textbf{Resolution} & \textbf{Gene coverage} & \textbf{Note} \\
\midrule
10x Visium & Sequencing & Spots $\sim$55\,\textmu m (a few cells each) & Whole transcriptome & Workhorse; spots are multi-cell, so not true single-cell. \\
Visium HD & Sequencing & 2--8\,\textmu m grid & Whole transcriptome & Near single-cell on a continuous grid. \\
Slide-seq & Sequencing & $\sim$10\,\textmu m beads & Whole transcriptome & Bead-array; high spatial detail, sparser per bead. \\
MERFISH & Imaging (ISH) & Subcellular & 100s--1000s targeted genes & Single-molecule, combinatorial barcodes. \\
seqFISH & Imaging (ISH) & Subcellular & 100s--1000s targeted genes & Sequential hybridization rounds. \\
In-situ sequencing & Imaging (ISS) & Subcellular & Targeted & Sequences short barcodes in place. \\
PhenoCycler / CODEX & Imaging (protein) & Subcellular & $\sim$40--100 proteins & Measures protein, not RNA --- complementary. \\
\bottomrule
\end{tabular}
\end{table}

\fbox{\parbox{\linewidth}{\small
\textbf{The one-line trade-off.} Sequencing-based methods (Visium, Slide-seq) win on gene breadth but are coarser
in space and often catch multiple cells per spot. Imaging-based methods (MERFISH, seqFISH) win on resolution and
sensitivity but only for a targeted panel. Choose your corner of the triangle to match the question.
}}

\section{What you can compute with spatial data}

\subsection{Reconstructing an axis $\rightarrow$ a zonation index}
Once every measurement has an $(x,y)$, expression becomes a function of space. If a tissue has a dominant geometric
axis (liver's porto-central axis, a cortical depth, a crypt--villus axis in gut), you can project each spot onto
that axis and ask how a gene varies along it.
\[
  \mathrm{zone\text{-}index}(\text{cell}) \;=\; \text{position of cell along the porto-central axis} \;\in\; [0,1].
\]
Concretely: pick a panel of genes known to be high at one pole and low at the other, combine them into a continuous
score, and you get a quantitative \emph{zonation index} --- a single number per cell saying ``how portal vs.\ how
central am I.'' This turns a vague anatomical notion into a measurable coordinate.

\subsection{Niche / neighborhood analysis}
For each cell, look at the cell types within a small radius and summarize that local mixture. Recurring mixtures
define \emph{niches} --- reproducible micro-environments (e.g.\ ``hepatocyte + endothelial + Kupffer''). You can
then test whether a niche changes in disease.

\subsection{Spatial ligand--receptor inference}
Signaling needs proximity. With coordinates you score ligand--receptor pairs only where the partners are actually
adjacent, distinguishing real local crosstalk from coincidence.

\section{Limitations --- keep these honest}
\fbox{\parbox{\linewidth}{\small
\textbf{Caveats.}
\begin{itemize}[leftmargin=1.4em, topsep=2pt]
  \item \textbf{Resolution vs.\ depth:} the more genes you read (sequencing-based), the coarser the spot; the finer
        the spot (imaging-based), the fewer genes --- you rarely get both.
  \item \textbf{Multi-cell spots:} a Visium spot mixes a few cells, so its profile is a blend. Recovering per-cell
        composition needs computational deconvolution against a single-cell reference.
  \item \textbf{Dropout \& sparsity:} lowly-expressed genes are missed (a zero may mean ``absent'' or merely ``not
        captured''), and per-spot/per-bead counts can be shallow.
  \item \textbf{2D slices of 3D tissue:} a section is one plane; axes and niches are reconstructed from a sample,
        not the whole organ.
\end{itemize}
}}

\section{Connecting to this project}
Here is the move at the heart of the hackathon, stated plainly.

\fbox{\parbox{\linewidth}{\small
\textbf{The big idea: spatial as a ruler.} A spatial dataset gives us a \emph{reference map} in which expression is
tied to position --- for the liver, where each expression profile sits along the porto-central zonation axis. We
learn, from that reference, the relationship ``this expression pattern $\rightarrow$ this zone-index.''
}}

Now take a dissociated dataset --- snRNA-seq of disease tissue, where position was destroyed at the dissociation
step. On its own it has no coordinates. But because the spatial reference taught us the expression-to-position
mapping, we can apply that mapping back onto the dissociated cells and decode an inferred position for each one.
\[
\begin{aligned}
  \text{spatial reference (learn)} \quad &\rightarrow\quad \text{expression} \Rightarrow \text{zone-index}\\
  \text{dissociated disease cells (apply)} \quad &\rightarrow\quad \text{recover lost position}
\end{aligned}
\]
The spatial map is the ruler; the dissociated data are the measurements we re-place onto it. snRNA-seq matters here
precisely because the disease tissue is frozen/fibrotic --- nuclei survive where whole cells would not. We trade the
irreversible loss of position for the ability to profile hard tissue, and then buy the position back computationally
using the spatial reference.

\fbox{\parbox{\linewidth}{\small
\textbf{Why this is powerful.} It lets us ask spatial questions --- ``is zonation disrupted in disease? which zone
do the sick cells come from?'' --- on datasets that have no spatial information of their own. The spatial reference
effectively donates its coordinate system to dissociated data.
}}

\subsection*{One-screen summary}
\begin{itemize}[leftmargin=1.4em]
  \item \textbf{Dissociated (sc/snRNA-seq):} all genes, no position.
  \item \textbf{Spatial (Visium, MERFISH, \ldots):} position kept, trade-offs in breadth/resolution.
  \item With position you can reconstruct axes (zonation), niches, and local signaling.
  \item \textbf{This project:} use the spatial axis as a ruler to decode lost positions onto dissociated disease
        cells.
\end{itemize}

\bigskip
\noindent\textit{\small Synthesis primer for the hackathon. Single-cell/sequencing foundations grounded in Genomics
Lecture~1; spatial-transcriptomics content compiled from standard methods knowledge and the project's framing, since
it is not in the typed lecture notes.}

\end{document}
